\documentclass[11pt,a4paper]{article}

% ── Packages ──────────────────────────────────────────────────────────────────
\usepackage[a4paper,top=1.9cm,bottom=2.0cm,left=2.0cm,right=2.0cm,headsep=10pt]{geometry}
\usepackage[T1]{fontenc}
\usepackage{lmodern}
\usepackage{microtype}
\usepackage{booktabs}
\usepackage{tabularx}
\usepackage{array}
\usepackage{enumitem}
\usepackage{xcolor}
\usepackage{listings}
\usepackage{titlesec}
\usepackage{fancyhdr}
\usepackage[hidelinks]{hyperref}
\usepackage{parskip}

% ── Spacing tightening ────────────────────────────────────────────────────────
\setlength{\parskip}{3pt plus 1pt minus 1pt}
\setlength{\textfloatsep}{6pt plus 2pt minus 2pt}
\setlength{\intextsep}{6pt plus 2pt minus 2pt}
\setlength{\abovecaptionskip}{3pt}
\setlength{\belowcaptionskip}{0pt}
\setlength{\tabcolsep}{4pt}
\renewcommand{\arraystretch}{1.05}
\setlist[itemize]{nosep,leftmargin=1.4em,topsep=2pt,partopsep=0pt}
\setlist[enumerate]{nosep,leftmargin=1.4em,topsep=2pt,partopsep=0pt}
\flushbottom

% ── Colours ───────────────────────────────────────────────────────────────────
\definecolor{kbdgray}{HTML}{F4F4F4}
\definecolor{rulekey}{HTML}{2A4D8F}
\definecolor{rulecomment}{HTML}{6C7A89}
\definecolor{accent}{HTML}{C0392B}

% ── Listings (code/rules) ─────────────────────────────────────────────────────
\lstdefinestyle{tinycode}{
    basicstyle=\ttfamily\footnotesize,
    backgroundcolor=\color{kbdgray},
    keywordstyle=\color{rulekey}\bfseries,
    commentstyle=\color{rulecomment}\itshape,
    stringstyle=\color{accent},
    breaklines=true,
    columns=fullflexible,
    keepspaces=true,
    showstringspaces=false,
    frame=single,
    rulecolor=\color{kbdgray},
    framerule=0pt,
    framesep=3pt,
    aboveskip=3pt,
    belowskip=3pt,
    xleftmargin=3pt,
    xrightmargin=3pt,
    morekeywords={match,allow,if,function,return,service,rules_version,async,await,final,const,if,else,Future,Stream,class,extends,void,bool,String,int,double,List,Map}
}
\lstset{style=tinycode}

% ── Section formatting ────────────────────────────────────────────────────────
\titleformat{\section}{\large\bfseries\color{rulekey}}{\thesection}{0.5em}{}
\titleformat{\subsection}{\normalsize\bfseries}{\thesubsection}{0.5em}{}
\titlespacing*{\section}{0pt}{0.7\baselineskip}{0.15\baselineskip}
\titlespacing*{\subsection}{0pt}{0.35\baselineskip}{0.05\baselineskip}

% ── Header ────────────────────────────────────────────────────────────────────
\pagestyle{fancy}
\fancyhf{}
\renewcommand{\headrulewidth}{0.4pt}
\fancyhead[L]{\small ClubConnect --- Enterprise Audit \& Orchestration Report}
\fancyhead[R]{\small Page \thepage}

% ── Title block is rendered manually below for tight, even spacing ───────────

\begin{document}
\thispagestyle{fancy}

\begingroup
\centering
{\LARGE\bfseries D2 Enterprise Audit \& Orchestration Report\par}
\vspace{2pt}
{\large ClubConnect --- Flutter community room application\par}
\vspace{2pt}
{\small 2026-05-29\par}
\endgroup
\vspace{4pt}
\hrule height 0.3pt
\vspace{6pt}

\noindent\textbf{Scope.} This report audits the ClubConnect Flutter application across four
dimensions: (1) the multi-agent orchestration workflow used during development;
(2) the domain architecture and Firestore data model; (3) the role-based access
control (RBAC) matrix and Firestore Security Rules; and (4) the observability
posture (Firebase Crashlytics) together with a forward-looking feature-flag
rollback plan. The report cites the repository at the commit on branch
\texttt{backend/google\_auth}. File paths in this document are repository-relative.

% ═══════════════════════════════════════════════════════════════════════════════
\section{Agent Workflow}
\label{sec:agent}

\subsection{Orchestration model}

Development of ClubConnect made systematic use of Claude Code as a multi-agent
orchestration platform. Two agent tiers were used: a \emph{primary} agent
handling synthesis, file edits, and human dialogue; and \emph{subagents}
(notably the read-only \texttt{Explore} agent) dispatched for codebase-wide
investigation. The primary agent never read the full \texttt{lib/} tree
directly; instead it delegated wide reads to the \texttt{Explore} subagent and
consumed a structured markdown report in return. This pattern keeps the primary
context window free of raw file noise, allowing longer multi-turn sessions
without compaction.

\begin{table}[h]
\centering\small
\begin{tabularx}{\linewidth}{@{}lXl@{}}
\toprule
\textbf{Agent} & \textbf{Responsibility} & \textbf{Tools used} \\
\midrule
Primary (Opus 4.7) & Synthesis, edits, user dialogue, plan & \texttt{Edit, Write, Bash} \\
\texttt{Explore} subagent & Read-only codebase survey & \texttt{Read, Grep, Glob} \\
\texttt{Plan} subagent & Architecture decisions (selective) & read-only \\
\bottomrule
\end{tabularx}
\caption{Agent tiers used during the build and audit phases.}
\end{table}

\subsection{Representative prompts}

Prompts were written as briefings rather than commands --- each subagent
prompt opens with the user-goal, the prior work already done, and the required
output shape. A representative invocation for the audit pass:

\begin{lstlisting}[language=,basicstyle=\ttfamily\scriptsize]
Agent({ description: "Audit models, providers, services",
        subagent_type: "Explore",
        prompt:
  "Inventory every model under lib/models/, every Firestore collection
   path touched in lib/services/, every provider in lib/providers/. For
   each provider list service deps and exposed state. Cross-check for
   firebase_crashlytics usages, feature flags, and analytics. Return
   structured markdown, 1500-2500 words. Cite file paths." })
\end{lstlisting}

Prompts that produced shallow output were rewritten to include (i) explicit
file lists, (ii) a target word count, and (iii) the consumer's downstream
need (``so I can cite this in a LaTeX security section'').

\subsection{Challenges: drift, mismatched references, and handoff cost}

Three recurring frictions were observed.

\textbf{Hallucinated file references.} A user instruction once referred to
\texttt{payment\_service.dart}, which does not exist; the actual payment
logic lives in \texttt{smart\_bill\_service.dart}. The orchestration policy
adopted was: \emph{verify before instrumenting}. The primary agent confirmed
the absence with \texttt{find lib -name "payment*"} before applying the
edit, and surfaced the mismatch in the response.

\textbf{Role-name collision.} The token \texttt{admin} denotes two distinct
concepts in this codebase --- a \emph{global} application role
(\texttt{users/\{uid\}.role}) and a \emph{community-scoped} role
(\texttt{communities/\{cid\}/members/\{uid\}.role}). Early prompts conflated
the two; later prompts and rule audits scope the word explicitly.

\textbf{Context drift between turns.} Long sessions risked summarisation
clipping load-bearing facts (e.g.\ ``the rules file does not cover
\texttt{events}''). The mitigation was to record critical findings into
durable artefacts (this report, in-code comments, structured commit
messages) rather than rely on transcript memory.

\subsection{Feature branch management}

The team used a feature-branch model on top of \texttt{main}. Recent active
branches included \texttt{backend/google\_auth} (current),
\texttt{backend/events}, and short-lived integration branches for the bill /
payment subsystem. Conventions in observed commit history:

\begin{itemize}[leftmargin=*,nosep]
  \item Feature work is scoped per subsystem (\texttt{backend/events},
        \texttt{backend/google\_auth}); cross-cutting concerns
        (\texttt{firebase\_crashlytics}) were applied on the active branch
        and merged forward.
  \item Agent-assisted edits are reviewed locally before commit; the
        primary agent never pushes or force-pushes without explicit
        user authorisation.
  \item Generated files (\texttt{lib/firebase\_options.dart}) are
        excluded from agent-driven edits per the working rules in
        \texttt{CLAUDE.md}; only \texttt{flutterfire configure}
        regenerates them.
\end{itemize}

% ═══════════════════════════════════════════════════════════════════════════════
\section{Architecture \& Data}
\label{sec:arch}

\subsection{Layered architecture}

ClubConnect follows a four-layer separation:
\textbf{UI} (\texttt{features/*}) $\rightarrow$
\textbf{Providers} (\texttt{providers/*}) $\rightarrow$
\textbf{Services} (\texttt{services/*}) $\rightarrow$
\textbf{Firebase / AI}.
No screen imports Firebase or the AI client directly; all I/O is routed
through services. Models are plain Dart (no Flutter imports), exposing
only \texttt{fromFirestore} / \texttt{toFirestore} round-trips.

\subsection{Domain model inventory}

Twenty-three model classes are defined in \texttt{lib/models/}. The
following grouping summarises the production-relevant entities:

\begin{table}[h]
\centering\small
\begin{tabularx}{\linewidth}{@{}lXl@{}}
\toprule
\textbf{Model} & \textbf{Purpose} & \textbf{Primary collection} \\
\midrule
\texttt{UserModel}          & Profile, role, ban / mute state          & \texttt{users/\{uid\}} \\
\texttt{CommunityModel}     & Community metadata, trending stats       & \texttt{communities/\{cid\}} \\
\texttt{MemberModel}        & Community membership + role              & \texttt{communities/\{cid\}/members} \\
\texttt{MessageModel}       & Chat message, mentions, reply threading  & \texttt{communities/\{cid\}/messages} \\
\texttt{EventModel}         & Event in a community                     & \texttt{communities/\{cid\}/events} \\
\texttt{AttendeeModel}      & Event attendance / host flag             & \texttt{.../events/\{eid\}/attendees} \\
\texttt{SmartBillModel}     & Event bill (split payments, AI verify)   & \texttt{.../events/\{eid\}/event\_bills} \\
\texttt{SmartBillItemModel} & Line item, per-payer share               & \texttt{.../event\_bills/\{bid\}/items} \\
\texttt{SmartPaymentModel}  & Per-user payment + AI verification       & \texttt{.../event\_bills/\{bid\}/payments} \\
\texttt{NotificationModel}  & User-scoped notification                 & \texttt{users/\{uid\}/notifications} \\
\texttt{ReviewModel}        & Numeric community review                 & \texttt{users/\{uid\}/rating} \\
\texttt{CategoryModel}      & Global tag registry                      & \texttt{categories/\{cid\}} \\
\texttt{ReportModel}        & Content / user report (admin queue)      & \texttt{reports/\{rid\}} \\
\texttt{RuleModel}          & Community rules (AI moderation input)    & embedded in community \\
\bottomrule
\end{tabularx}
\caption{Domain models and the Firestore documents they project to.}
\end{table}

\subsection{Sub-collection hierarchy}

The Firestore tree is intentionally nested so that authorisation can be
derived from the document path. Full paths in use:

\begin{lstlisting}[basicstyle=\ttfamily\scriptsize]
users/{uid}
users/{uid}/notifications/{notificationId}
users/{uid}/rating/{ratingId}
users/{uid}/interestScores/{tagSlug}

communities/{communityId}
communities/{communityId}/members/{userId}
communities/{communityId}/messages/{messageId}
communities/{communityId}/events/{eventId}
communities/{communityId}/events/{eventId}/attendees/{userId}
communities/{communityId}/events/{eventId}/messages/{messageId}
communities/{communityId}/events/{eventId}/event_bills/{billId}
communities/{communityId}/events/{eventId}/event_bills/{billId}/items/{itemId}
communities/{communityId}/events/{eventId}/event_bills/{billId}/payments/{userId}

community_activity/{communityId}/activeUsers/{userId}
categories/{categoryId}
reports/{reportId}
\end{lstlisting}

The deepest live path is six levels deep (a payment under a bill under an
event under a community). This is intentional: it lets a single security
predicate at the \texttt{community} level transitively gate every nested
document via membership / role checks.

\subsection{State management}

The team chose the \texttt{provider} package (\texttt{ChangeNotifier}) over
Riverpod or BLoC. The justification is three-fold:

\begin{enumerate}[leftmargin=*,nosep]
  \item \textbf{Ceremony cost.} \texttt{ChangeNotifier} requires no code
        generation and no learning curve beyond \texttt{setState}. For a
        capstone team, build velocity outranked compile-time guarantees.
  \item \textbf{Stream-friendly.} Each provider holds named
        \texttt{StreamSubscription} handles to Firestore queries
        (\texttt{\_billSub}, \texttt{\_itemsSub}, \texttt{\_paymentSub} in
        \texttt{SmartBillProvider}) and cancels them in
        \texttt{dispose()}, eliminating the listener-leak class of bugs.
  \item \textbf{Boundary discipline.} The strict layer rule (\emph{screens
        never call services}) is naturally enforced by exposing only
        provider methods on UI widgets.
\end{enumerate}

Ten providers are wired at the root in \texttt{lib/main.dart}, one per
domain (\texttt{AppAuthProvider}, \texttt{CommunityProvider},
\texttt{EventProvider}, \texttt{SmartBillProvider}, etc.). Cross-provider
coordination is achieved by injection: e.g.\ \texttt{SmartBillProvider}
takes a \texttt{NotificationService} constructor argument so that a
verified payment can write a notification to the host.

\subsection{Key design decisions}

\begin{description}[leftmargin=*,style=nextline,nosep]
  \item[Route argument wrappers.] Complex inter-route data
        (\texttt{EventDetailArgs}, \texttt{ChatArgs},
        \texttt{SmartPayBillArgs}) is passed via typed wrapper classes
        through \texttt{GoRouter}'s \texttt{extra} channel rather than raw
        maps, keeping the router contract type-safe.
  \item[Pre-write AI moderation.] \texttt{GeminiService.moderateMessage()}
        is invoked \emph{before} any Firestore write in
        \texttt{chat\_screen.dart}; blocked content never persists. This
        avoids retroactive cleanup, which is expensive on a nested
        schema.
  \item[Event-scoped chat rooms.] \texttt{RoomModel} distinguishes
        community rooms (permanent) from event rooms (\texttt{expiresAt}
        bound to event end). Expiry is a query filter, not a delete.
  \item[Lazy stats windows.] \texttt{CommunityStats} carries a 24-hour
        rolling window with a lazy reset (\texttt{statsWindowStart}),
        avoiding scheduled Cloud Functions for the trending feed.
\end{description}

% ═══════════════════════════════════════════════════════════════════════════════
\section{Security Matrix}
\label{sec:sec}

\subsection{RBAC matrix}

ClubConnect recognises three role tokens in two namespaces. Global roles
live on the user document; community roles live on each membership record.

\begin{table}[h]
\centering\small
\renewcommand{\arraystretch}{1.15}
\begin{tabularx}{\linewidth}{@{}lXXX@{}}
\toprule
\textbf{Capability} & \textbf{user (default)} & \textbf{admin (global)} & \textbf{admin / creator (community)} \\
\midrule
Read user profiles                 & Yes & Yes & Yes \\
Edit own profile                   & Yes & Yes & Yes \\
Edit another user's profile        & No  & No  & No \\
Browse communities                 & Yes & Yes & Yes \\
Create community                   & Yes & Yes & Yes \\
Edit community metadata            & No  & No  & Yes \\
Add / remove members               & Self only & Self only & Yes \\
Change member roles                & No  & No  & Yes (creator only for admin promotion) \\
Post / read messages               & If member & If member & If member \\
Delete others' messages            & No  & No  & Yes (admin / creator of community) \\
Submit report                      & Yes & Yes & Yes \\
Review / resolve reports           & No  & Yes (server-side) & No \\
Ban / mute user                    & No  & Yes & No \\
Access \texttt{/admin} routes      & No  & Yes & No \\
\bottomrule
\end{tabularx}
\caption{Role $\times$ capability matrix. ``Server-side'' indicates
operations that the client cannot perform directly --- the security rules
forbid the write, and the admin UI must invoke a privileged path (Cloud
Function or trusted service account).}
\end{table}

The global \texttt{admin} check is enforced at two points: the router
redirect in \texttt{lib/router/app\_router.dart} blocks navigation to
\texttt{/admin*} routes if \texttt{authProvider.role != 'admin'}, and the
\texttt{AdminReportsScreen} guard re-checks the role on mount as a
defence-in-depth measure.

\subsection{Firestore Security Rules}

The deployed rule file (\texttt{firestore.rules}) declares five reusable
predicates and four explicit resource matches before falling through to a
deny-everything catch-all.

\begin{lstlisting}[basicstyle=\ttfamily\scriptsize]
function isMember(communityId) {
  return isAuthed() &&
    exists(/databases/$(database)/documents/communities/$(communityId)
           /members/$(request.auth.uid));
}
function isAdminOrCreator(communityId) {
  return isMember(communityId) &&
    memberRole(communityId) in ['admin', 'creator'];
}
\end{lstlisting}

The rules are summarised below; predicates in parentheses are the
constraint enforced by the rule.

\begin{table}[h]
\centering\small
\renewcommand{\arraystretch}{1.1}
\begin{tabularx}{\linewidth}{@{}lXX@{}}
\toprule
\textbf{Path} & \textbf{Read} & \textbf{Write} \\
\midrule
\texttt{users/\{uid\}}                          & \texttt{isAuthed()} & \texttt{create/update}: \texttt{isCurrentUser}; \texttt{delete}: deny \\
\texttt{users/\{uid\}/notifications/\{n\}}      & owner only           & owner only \\
\texttt{communities/\{cid\}}                    & \texttt{isAuthed()}  & \texttt{create}: any authed; \texttt{update}: \texttt{isAdminOrCreator}; \texttt{delete}: deny \\
\texttt{.../members/\{uid\}}                    & \texttt{isAuthed()}  & \texttt{create}: self or admin / creator; \texttt{update}: admin / creator; \texttt{delete}: self or admin / creator \\
\texttt{.../messages/\{mid\}}                   & \texttt{isMember}    & \texttt{create}: \texttt{isMember}; \texttt{update}: sender, or \texttt{seenBy}-only diff; \texttt{delete}: sender or admin / creator \\
\texttt{reports/\{rid\}}                        & \texttt{isAuthed()}  & \texttt{create}: any authed; \texttt{update/delete}: deny \\
\texttt{/\{document=**\}}                       & deny                 & deny \\
\bottomrule
\end{tabularx}
\caption{Effective rule matrix as written in \texttt{firestore.rules}.}
\end{table}

The \texttt{messages.update} rule is particularly defensive: a non-sender
may update a message \emph{only} if the diff affects exclusively the
\texttt{seenBy} field. This permits read-receipts without enabling content
tampering.

\subsection{Findings and gaps}

The repository's \texttt{firestore.rules} explicitly grants access to
\texttt{users}, \texttt{users/notifications}, \texttt{communities},
\texttt{communities/members}, \texttt{communities/messages}, and
\texttt{reports}. Because the rule set ends with
\texttt{match /\{document=**\} \{ allow read, write: if false; \}}, every
unlisted path is denied.

The following live application paths are \emph{not} covered by any
\texttt{allow} rule and would currently be blocked by the catch-all if the
rules in source were the deployed rules:

\begin{itemize}[leftmargin=*,nosep]
  \item \texttt{communities/\{cid\}/events/\{eid\}} and all nested
        subcollections (\texttt{attendees}, \texttt{messages},
        \texttt{event\_bills}, \texttt{items}, \texttt{payments}).
  \item \texttt{community\_activity/\{cid\}/activeUsers/\{uid\}}.
  \item \texttt{categories/\{cid\}}.
  \item \texttt{users/\{uid\}/rating/\{rid\}} (consumed by
        \texttt{RatingProvider} and \texttt{ReviewModel}).
\end{itemize}

\textbf{Remediation.} Add explicit \texttt{match} blocks for each path
above, gated on \texttt{isMember(communityId)} for community-scoped paths
and \texttt{isCurrentUser(userId)} for user-scoped paths, before the next
deployment. Until this is done, the deployed rule set must remain a
superset (typically a console-edited permissive policy) or production will
break.

% ═══════════════════════════════════════════════════════════════════════════════
\section{Observability \& Rollback}
\label{sec:obs}

\subsection{Crashlytics integration map}

Firebase Crashlytics has been added to \texttt{pubspec.yaml}
(\texttt{firebase\_crashlytics: \^{}4.0.0}, resolved to 4.3.10) and wired
at five sites. Suppressed-error policy: \emph{expected} authentication
failures (\texttt{wrong-password}, \texttt{user-not-found},
\texttt{invalid-email}, \texttt{invalid-credential}) are \emph{not}
recorded --- only unexpected codes are.

\begin{table}[h]
\centering\small
\begin{tabularx}{\linewidth}{@{}lXl@{}}
\toprule
\textbf{File} & \textbf{Coverage} & \textbf{Severity} \\
\midrule
\texttt{lib/main.dart}                       & Zone-level catch via \texttt{runZonedGuarded}; \texttt{FlutterError.onError}; \texttt{PlatformDispatcher.instance.onError}; collection disabled in debug & fatal \\
\texttt{lib/services/notification\_service.dart} & try / catch on every write (\texttt{createNotification}, \texttt{markAsRead}, \texttt{markAllAsRead}, \texttt{deleteNotification}, \texttt{getNotifications}) plus stream \texttt{handleError} on \texttt{watchNotifications} & non-fatal \\
\texttt{lib/providers/smart\_bill\_provider.dart} & \texttt{\_recordError} helper invoked in every catch (\texttt{deleteBill}, \texttt{settleBill}, \texttt{verifyMemberPayment}, \texttt{submitAndVerify}, \texttt{submitPaymentWithSlip}, \texttt{verifySlip}, \texttt{create / updateAndPublishBill}) & non-fatal \\
\texttt{lib/services/smart\_bill\_service.dart} & try / catch around \texttt{uploadQrImage}, \texttt{uploadSlip}, \texttt{verifySlip} (AI call) & non-fatal \\
\texttt{lib/providers/auth\_provider.dart} & \texttt{sendOtp}, \texttt{verifyOtp}, \texttt{signUp}, \texttt{startGoogleRegistration}, user-doc stream \texttt{onError}; \texttt{signIn} suppresses expected codes & non-fatal \\
\bottomrule
\end{tabularx}
\caption{Crashlytics instrumentation sites and severity.}
\end{table}

\subsection{Initialization}

The application entry point in \texttt{lib/main.dart} wraps the runtime in
a guarded zone, attaches the framework- and platform-level error
forwarders, and disables collection in debug builds:

\begin{lstlisting}[basicstyle=\ttfamily\scriptsize]
void main() {
  runZonedGuarded<Future<void>>(() async {
    WidgetsFlutterBinding.ensureInitialized();
    await dotenv.load(fileName: '.env');
    await Firebase.initializeApp(
        options: DefaultFirebaseOptions.currentPlatform);

    await FirebaseCrashlytics.instance
        .setCrashlyticsCollectionEnabled(!kDebugMode);

    FlutterError.onError = FirebaseCrashlytics.instance.recordFlutterFatalError;
    PlatformDispatcher.instance.onError = (error, stack) {
      FirebaseCrashlytics.instance.recordError(error, stack, fatal: true);
      return true;
    };

    runApp(const ClubConnectApp());
  }, (error, stack) {
    FirebaseCrashlytics.instance.recordError(error, stack, fatal: true);
  });
}
\end{lstlisting}

Per-call records include a \texttt{reason} string and a contextual
\texttt{information} payload (community / event / bill / user IDs) so that
incidents can be triaged without raw stack archaeology.

\subsection{Feature flags --- current state and proposed plan}

\textbf{Current state.} ClubConnect has \emph{no} runtime feature-flag
mechanism. There are no usages of \texttt{firebase\_remote\_config}, no
kill switches, and no environment-driven gating beyond the
\texttt{kDebugMode} check that scopes Crashlytics collection. Rollback
today is exclusively a code-revert-and-redeploy operation.

\textbf{Proposed plan.} Introduce Firebase Remote Config with a small set
of flags, each gating an externally dependent or experimental subsystem.
Suggested flag set:

\begin{table}[h]
\centering\small
\renewcommand{\arraystretch}{1.1}
\begin{tabularx}{\linewidth}{@{}lXl@{}}
\toprule
\textbf{Flag} & \textbf{Gates} & \textbf{Default} \\
\midrule
\texttt{ai\_moderation\_enabled}         & Pre-send AI moderation in \texttt{chat\_screen.dart}      & true  \\
\texttt{ai\_bill\_verification\_enabled} & AI slip verification in \texttt{SmartBillService.verifySlip} & true  \\
\texttt{google\_auth\_enabled}           & Google sign-in button on welcome / signup screens         & true  \\
\texttt{phone\_otp\_enabled}             & Phone OTP path in signup flow                              & true  \\
\texttt{trending\_feed\_enabled}         & Trending-communities computation on home tab               & true  \\
\bottomrule
\end{tabularx}
\caption{Proposed Remote Config flags. All default to \texttt{true}; an
incident toggles to \texttt{false} as the first rollback step.}
\end{table}

\subsection{Rollback plan and reversibility tiers}

Operational rollback should escalate through three reversibility tiers,
matching the \texttt{CLAUDE.md} convention:

\begin{enumerate}[leftmargin=*,nosep]
  \item \textbf{R2 (easily reversed) --- Remote Config flip.} Set the
        offending flag to \texttt{false} in the Firebase console.
        Propagation time is bounded by the client fetch interval (default
        12 hours; recommend setting \texttt{minimumFetchInterval} to 1
        hour for fast incident response). No deploy required.
  \item \textbf{R1 (costly to reverse) --- Tighten Firestore rules.} If
        a write path is the root cause (e.g.\ a malformed bill document
        type), publish a rule revision that denies the offending write
        shape. This is fast (seconds) but requires a rule audit afterwards
        because the change is global.
  \item \textbf{R0 (irreversible without effort) --- App rollback.}
        Re-release a prior signed build via the Play Store /
        TestFlight rollout channel. Requires user reinstall or
        background update; only used when client logic is irrecoverably
        broken.
\end{enumerate}

\textbf{Crashlytics as the rollback trigger.} The Crashlytics dashboard is
the primary signal: a crash-free-users-rate drop below 99.0\% or any
non-fatal recorded under \texttt{SmartBillProvider.submitAndVerify} (a
revenue-adjacent path) should escalate to a Remote Config flip within
fifteen minutes. The \texttt{information} fields recorded at each
instrumentation site (community / event / bill IDs) are sufficient to
scope the flip to a specific cohort without disabling the feature
globally if a future flag scheme supports targeting.

% ═══════════════════════════════════════════════════════════════════════════════
\section*{Appendix A --- File inventory}

\noindent\textbf{Touched by the Crashlytics integration:} \texttt{pubspec.yaml},
\texttt{lib/main.dart}, \texttt{lib/services/notification\_service.dart},
\texttt{lib/services/smart\_bill\_service.dart},
\texttt{lib/providers/smart\_bill\_provider.dart},
\texttt{lib/providers/auth\_provider.dart}.

\noindent\textbf{Touched by the data hardening pass:}
\texttt{lib/models/review\_model.dart} (numeric coercion fix:
\texttt{(json['score'] as num?)?.toDouble()} replaces the unsafe
\texttt{as double} cast that crashed on integer Firestore values).

\noindent\textbf{Not yet touched (remediation candidates):}
\texttt{firestore.rules} (missing event / bill / payment / category /
review paths; see \S\ref{sec:sec}); \texttt{pubspec.yaml} (no
\texttt{firebase\_remote\_config}; see \S\ref{sec:obs});
\texttt{android/app/build.gradle} (Crashlytics Gradle plugin for native
symbol upload is not yet wired --- Dart-side reporting works without it,
but native NDK symbolication would not).

\vspace{0.6em}\noindent\rule{\linewidth}{0.3pt}\\
\small \emph{End of report. Compiled from repository state on
branch \texttt{backend/google\_auth} as of 29 May 2026.}

\end{document}
